\chapter{Implementation}

This chapter presents the practical implementation details of the Binary Search and Jump Search algorithms within the Dataclenz library, highlighting key code snippets and explaining implementation decisions.

\section{Binary Search Implementation}

The Binary Search algorithm is implemented as follows:

\begin{lstlisting}[style=c]
int binary_search(DataFrame* df, int column_index, void* target, int* found_indices, int max_indices) {
    if (df == NULL || column_index < 0 || column_index >= df->num_columns || found_indices == NULL) {
        set_error("Invalid parameters for binary search");
        return 0;
    }

    Column* column = &df->columns[column_index];
    int found_count = 0;
    
    // Check if array is empty
    if (df->num_rows == 0) {
        return 0;
    }

    int left = 0;
    int right = df->num_rows - 1;
    int mid = 0;
    int compare_result = 0;

    // Binary search to find any occurrence
    while (left <= right) {
        mid = left + (right - left) / 2;
        
        // Compare based on data type
        switch (column->type) {
            case TYPE_INT:
                compare_result = *(int*)target - ((int*)column->data)[mid];
                break;
            case TYPE_FLOAT:
                {
                    float diff = *(float*)target - ((float*)column->data)[mid];
                    compare_result = (diff > 0) ? 1 : ((diff < 0) ? -1 : 0);
                }
                break;
            case TYPE_STRING:
                compare_result = strcmp((char*)target, ((char**)column->data)[mid]);
                break;
            default:
                set_error("Unsupported data type for binary search");
                return 0;
        }

        if (compare_result == 0) {
            // Found a match
            found_indices[found_count++] = mid;
            
            // Check for more matches to the left and right
            // ... [code omitted for brevity] ...
            
            return found_count;
        } else if (compare_result < 0) {
            right = mid - 1;
        } else {
            left = mid + 1;
        }
    }
    
    return found_count;
}
\end{lstlisting}

\section{Jump Search Implementation}

The Jump Search algorithm is implemented as follows:

\begin{lstlisting}[style=c]
int jump_search(DataFrame* df, int column_index, void* target, int* found_indices, int max_indices) {
    if (df == NULL || column_index < 0 || column_index >= df->num_columns || found_indices == NULL) {
        set_error("Invalid parameters for jump search");
        return 0;
    }

    Column* column = &df->columns[column_index];
    int found_count = 0;
    int n = df->num_rows;
    
    if (n == 0) {
        return 0;
    }
    
    // Finding block size to be jumped
    int step = (int)sqrt(n);
    
    // Finding the block where element is present
    int prev = 0;
    int compare_result = 0;
    
    // Jump steps
    while (prev < n) {
        // Compare current element with target
        switch (column->type) {
            case TYPE_INT:
                compare_result = *(int*)target - ((int*)column->data)[prev];
                break;
            case TYPE_FLOAT:
                {
                    float diff = *(float*)target - ((float*)column->data)[prev];
                    compare_result = (diff > 0) ? 1 : ((diff < 0) ? -1 : 0);
                }
                break;
            case TYPE_STRING:
                compare_result = strcmp((char*)target, ((char**)column->data)[prev]);
                break;
            default:
                set_error("Unsupported data type for jump search");
                return 0;
        }
        
        if (compare_result == 0) {
            // Found a match
            found_indices[found_count++] = prev;
            
            // Check for more matches in vicinity
            // ... [code omitted for brevity] ...
            
            return found_count;
        } else if (compare_result < 0) {
            // Element is in this block, perform linear search
            break;
        }
        
        prev = (prev + step < n) ? prev + step : n;
    }
    
    // Linear search in identified block
    // ... [code omitted for brevity] ...
    
    return found_count;
}
\end{lstlisting}

\section{Performance Measurement Implementation}

To accurately measure algorithm performance, a specialized function was implemented:

\begin{lstlisting}[style=c]
SearchPerformanceMetrics measure_binary_search_performance(DataFrame* df, int column_index, int num_searches, ColumnType type) {
    SearchPerformanceMetrics metrics;
    strncpy(metrics.search_type, "Binary Search", sizeof(metrics.search_type) - 1);
    metrics.dataset_size = df->num_rows;
    metrics.num_searches = num_searches;
    metrics.avg_search_time = 0.0;
    metrics.avg_comparisons = 0;
    metrics.successful_searches = 0;
    metrics.memory_usage_kb = sizeof(int) * df->num_rows / 1024;

    // Sort the DataFrame first (binary search requires sorted data)
    DataFrame* sorted_df = sort_dataframe(df, column_index, 1);
    if (sorted_df == NULL) {
        metrics.avg_search_time = -1.0;  // Indicate failure
        return metrics;
    }

    // Allocate memory for found indices
    int* found_indices = (int*)malloc(df->num_rows * sizeof(int));
    if (found_indices == NULL) {
        metrics.avg_search_time = -1.0;  // Indicate failure
        return metrics;
    }

    // Measure search performance
    clock_t start_time = clock();
    int total_comparisons = 0;

    for (int i = 0; i < num_searches; i++) {
        // Generate random search target
        int random_index = rand() % df->num_rows;
        void* target;
        
        // Prepare target based on data type
        // ... [code omitted for brevity] ...

        // Perform binary search
        int found_count = binary_search(sorted_df, column_index, target, found_indices, df->num_rows);
        
        if (found_count > 0) {
            metrics.successful_searches++;
        }
        
        // Estimate comparisons (log2(n) for binary search)
        total_comparisons += (int)ceil(log2(df->num_rows));
        
        free(target);
    }

    clock_t end_time = clock();
    metrics.avg_search_time = ((double)(end_time - start_time)) / CLOCKS_PER_SEC / num_searches;
    metrics.avg_comparisons = total_comparisons / num_searches;

    free(found_indices);
    return metrics;
}
\end{lstlisting}

A similar function was implemented for Jump Search, with the primary difference being the estimation of comparisons (sqrt(n) for Jump Search).

\section{Key Implementation Decisions}

\subsection{Type-Agnostic Comparison}
To support multiple data types, the implementation uses a switch statement to perform the appropriate comparison based on the column type. This approach enables the search algorithms to work with integers, floats, and strings without requiring separate implementations for each type.

\subsection{Duplicate Handling}
Both search algorithms include mechanisms to handle duplicate values by scanning adjacent elements when a match is found. This ensures all occurrences of the target value are identified, which is crucial for real-world datasets where duplicates are common.

\subsection{Sort Requirement}
Since both Binary Search and Jump Search require sorted data, the implementation includes explicit sorting before searching. The sort operation is performed outside the timed section to focus measurement on search performance rather than preprocessing time.

\subsection{Memory Management}
The implementation carefully manages memory allocation and deallocation to prevent leaks. Temporary buffers for search results are allocated once per test and reused across multiple searches to minimize overhead.

\subsection{Error Handling}
Robust error checking is implemented throughout the code to ensure valid inputs and proper error reporting. This includes verification of DataFrame pointers, column indices, and allocated memory.

\section{Performance Optimization Techniques}

Several optimization techniques were employed to enhance algorithm performance:

\begin{itemize}
    \item \textbf{Mid-point Calculation:} Using the formula \texttt{mid = left + (right - left) / 2} to avoid potential integer overflow
    \item \textbf{Early Termination:} Exiting loops as soon as a match is found to minimize unnecessary comparisons
    \item \textbf{Memory Locality:} Accessing data sequentially where possible to improve cache utilization
    \item \textbf{Optimized Step Size:} Using sqrt(n) as the step size for Jump Search to balance jump steps and linear search effort
\end{itemize}

These optimizations ensure the implementation achieves performance close to the theoretical time complexity while remaining robust and reliable.